\chapter{Mã nguồn và khả năng tái lập}
\label{app:source-code}

\noindent
Mã nguồn của đồ án được công bố tại (thư mục \texttt{pinn-soh}):
\begin{center}
  \url{https://github.com/huynhquocthang1805/-n-cn}
\end{center}
Mã của GBDT lịch sử dùng trong E33 (\Cref{sec:e33pp}) lấy từ phần cải tiến (thư mục \texttt{improvements}) của repo
\begin{center}
  \url{https://github.com/huynhquocthang1805/G-doan}
\end{center}
và đặt trong \texttt{pinnsoh/forecast\_hist.py}. Dữ liệu là bản đặc trưng công bố kèm PINN4SOH tại
\url{https://github.com/wang-fujin/PINN4SOH}.

\section{Chạy lại}

Cần Python 3.11 và các thư viện trong \texttt{requirements.txt}; bước dựng báo cáo cần pdf\LaTeX{} và Bib\TeX.
\begin{center}\small
\begin{tabular}{@{}ll@{}}
\toprule
Lệnh & Việc làm \\
\midrule
\texttt{python run.py data} & tải dữ liệu, dựng bộ nhớ đệm \\
\texttt{python run.py E13} (\ldots\ \texttt{E33}) & chạy một thí nghiệm, bỏ qua lượt đã có kết quả \\
\texttt{python run.py analyze} & bảng và hình từ kết quả \\
\texttt{python run.py benchmark} & chạy lại mã gốc PINN4SOH, lập bảng so sánh \\
\texttt{python run.py features} & phân tích đặc trưng đầu vào \\
\texttt{python run.py report} & dựng báo cáo này \\
\texttt{python run.py reproduce} & chạy lại 20 lượt mẫu, so từng bit \\
\bottomrule
\end{tabular}
\end{center}

\section{Các kiểm tra tái lập}\label{sec:tailap}

Mọi lượt chạy dùng một luồng và hạt giống cố định, nên trên cùng loại CPU và cùng phiên bản thư viện, chạy lại phải cho dự đoán trùng từng bit. Các kiểm tra đã làm:

\begin{enumerate}
\item Hai mươi lượt mẫu phủ mọi loại thí nghiệm chính (MLP, PINN-semi, PINN-sup, PINN hộp đen, Ridge, GPR, GBDT, kiểm toán rò rỉ, chuyển miền), chạy lại vào thư mục tạm và so với giá trị lưu của lần chạy báo cáo đã nộp: 20 trên 20 lượt có MAE theo cell và mã băm của toàn bộ dự đoán test trùng từng bit.
\item Kiểm toán rò rỉ (160 lượt của bốn cách làm sai và 40 lượt tham chiếu theo cell) được chạy lại toàn bộ; 5 trên 5 lượt có giá trị tham chiếu lưu sẵn trùng từng bit, và các tỉ số ở Bảng~\ref{tab:leak} khớp với báo cáo đã nộp.
\item PINN-lai, GBDT và PINN-semi trên 39 nhóm của PINN4SOH (1.002 lượt) được chạy lại toàn bộ và so với kết quả của lần chạy trước: 216 trên 216 giá trị trung bình (MAE, RMSE và độ lệch chuẩn qua ba hạt giống của PINN-lai, GBDT và trung bình ba PINN-semi, với hai cách tiền xử lý, trên 12 chế độ) trùng từng bit; MAE và RMSE theo cell của cả 3.852 dòng trùng ở 6 chữ số có nghĩa đã lưu.
\item GBDT với đủ 17 đầu vào trong phân tích đặc trưng (Mục~\ref{sec:dt_pp}) cho MAE theo cell trùng từng bit với lượt GBDT tham chiếu của thí nghiệm chính (XJTU, 30\,\% nhãn, lần lặp 0, fold 0).
\item So với PINN4SOH trên đúng dữ liệu train của bài: một đại lượng ghi theo từng epoch trong nhật ký huấn luyện công bố của bài, tính lại trên các hàng train dựng lại ở đây, khớp ở cả 39 nhóm; sai số khi chạy lại mã gốc không trùng số công bố vì khác phần cứng và phiên bản thư viện.
\end{enumerate}

Khác CPU (ví dụ AVX-512 so với AVX2) hoặc khác phiên bản PyTorch có thể làm kết quả lệch ở mức làm tròn số thực; khi đó các kết luận không đổi nhưng chữ số cuối trong bảng có thể khác.

